% GENERATED FILE. Edit canonical lessons, then run npm run generate:books.
% canonical-source-hash: fnv1a64:d1cb2906bc75d760
% canonical-lessons: GU-C520-vagharvu, GU-C520-satalvu, GU-C520-capti, GU-C520-chatvu, GU-C520-bhujvu

\chapter{Spices in Hot Oil}
\label{ch:gu-spices-in-hot-oil}
\hlchaptermodality{520}

\begin{chapteropening}
\textbf{By the end of this chapter:} \emph{I can describe tempering, frying and roasting spices.}

Complete the last lesson of chapter 520: I can describe tempering, frying and roasting spices.
\end{chapteropening}

\section[\texorpdfstring{\gu{વઘારવું} (vaghārvũ)}{વઘારવું (vaghārvũ)}]{\gu{વઘારવું} (vaghārvũ) — to temper (with hot spiced oil)}

\label{lesson:GU-C520-vagharvu}



\begin{quote}
Before the new one: say the Gujarati for steam, then the Gujarati for to let cool.
\end{quote}

\subsection*{You'll want to know: \gu{વઘારવું}}
\textbf{\gu{વઘારવું}} — \emph{vaghārvũ} — \textquotedblleft{}to temper (with hot spiced oil)\textquotedblright{}.

\begin{quote}
\gu{માતાએ} \gu{રાઈ} \gu{અને} \gu{જીરાથી} \gu{દાળ} \gu{વઘારી}.
\end{quote}

\begin{quote}
\emph{mātāe rāī ane jīrāthī dāḷ vaghārī.}
\end{quote}

\begin{quote}
Mother tempered the dal with mustard seed and cumin.
\end{quote}

Past, with \gu{તેણે} (he or she …): \gu{તેણે} \gu{દાળ} \gu{વઘારી}, \emph{teṇe dāḷ vaghārī}.

\begin{cousinweb}
\textbf{\gu{વઘારવું}} (\emph{vaghārvũ}) is inherited from Sanskrit \emph{vyághārayati}, \textquotedblleft{}sprinkles over\textquotedblright{}, and CDIAL lists the Gujarati verb, \textquotedblleft{}to season\textquotedblright{} (CDIAL 12191): hot oil with spices is poured over the dish.
\end{cousinweb}

\begin{grammarlens}[title={What you've built}]
One more word for describing what you lived and felt.
\end{grammarlens}

\subsection*{Your turn}
\begin{itemize}
\raggedright
  \item \emph{Say it:} \emph{vaghārvũ}
  \item \emph{Say it:} \emph{vaghārvũ}, once more
  \item \emph{Say it:} \emph{mātāe rāī ane jīrāthī dāḷ vaghārī.} — Mother tempered the dal with mustard seed and cumin.
  \item \emph{Recall:} say the Gujarati for steam, then the Gujarati for to let cool, then say \emph{vaghārvũ} again
\end{itemize}

\begin{tcolorbox}[breakable,colback=teal!4,colframe=teal!35!black,title={Before you move on}]
What does \gu{વઘારવું} mean? (\textquotedblleft{}To temper (with hot spiced oil)\textquotedblright{}.) Say it once more.
\end{tcolorbox}
\section[\texorpdfstring{\gu{સાંતળવું} (sā̃taḷvũ)}{સાંતળવું (sā̃taḷvũ)}]{\gu{સાંતળવું} (sā̃taḷvũ) — to sauté}

\label{lesson:GU-C520-satalvu}



\begin{quote}
Before the new one: say the Gujarati for to let cool, then the Gujarati for to temper.
\end{quote}

\subsection*{You'll want to know: \gu{સાંતળવું}}
\textbf{\gu{સાંતળવું}} — \emph{sā̃taḷvũ} — \textquotedblleft{}to sauté\textquotedblright{}.

\begin{quote}
\gu{મીરાએ} \gu{તેલમાં} \gu{ડુંગળી} \gu{સાંતળી}.
\end{quote}

\begin{quote}
\emph{Mīrāe telmā̃ ḍungḷī sā̃tḷī.}
\end{quote}

\begin{quote}
Mira sautéed the onion in oil.
\end{quote}

Past, with \gu{તેણે} (he or she …): \gu{તેણે} \gu{ડુંગળી} \gu{સાંતળી}, \emph{teṇe ḍungḷī sā̃tḷī}.

\begin{cousinweb}
\textbf{\gu{સાંતળવું}} (\emph{sā̃taḷvũ}) comes from \emph{*saṁtalati}, a form CDIAL reconstructs, and CDIAL lists the Gujarati verb, \textquotedblleft{}to fry in ghee or oil\textquotedblright{} (CDIAL 12883).
\end{cousinweb}

\begin{grammarlens}[title={What you've built}]
One more word for describing what you lived and felt.
\end{grammarlens}

\subsection*{Your turn}
\begin{itemize}
\raggedright
  \item \emph{Say it:} \emph{sā̃taḷvũ}
  \item \emph{Say it:} \emph{sā̃taḷvũ}, once more
  \item \emph{Say it:} \emph{Mīrāe telmā̃ ḍungḷī sā̃tḷī.} — Mira sautéed the onion in oil.
  \item \emph{Recall:} say the Gujarati for to let cool, then the Gujarati for to temper, then say \emph{sā̃taḷvũ} again
\end{itemize}

\begin{tcolorbox}[breakable,colback=teal!4,colframe=teal!35!black,title={Before you move on}]
What does \gu{સાંતળવું} mean? (\textquotedblleft{}To sauté\textquotedblright{}.) Say it once more.
\end{tcolorbox}
\section[\texorpdfstring{\gu{ચપટી} (capṭī)}{ચપટી (capṭī)}]{\gu{ચપટી} (capṭī) — a pinch (of something)}

\label{lesson:GU-C520-capti}



\begin{quote}
Before the new one: say the Gujarati for to temper, then the Gujarati for to sauté.
\end{quote}

\subsection*{You'll want to know: \gu{ચપટી}}
\textbf{\gu{ચપટી}} — \emph{capṭī} — \textquotedblleft{}a pinch (of something)\textquotedblright{}.

\begin{quote}
\gu{શાકમાં} \gu{એક} \gu{ચપટી} \gu{હળદર} \gu{નાખો}.
\end{quote}

\begin{quote}
\emph{śākmā̃ ek capṭī haḷdar nākho.}
\end{quote}

\begin{quote}
Put a pinch of turmeric in the vegetables.
\end{quote}

A feminine noun.

\begin{cousinweb}
\textbf{\gu{ચપટી}} (\emph{capṭī}) is as much as fits between finger and thumb. No dictionary we use records an older form, so its history is not settled.
\end{cousinweb}

\begin{grammarlens}[title={What you've built}]
One more word for describing what you lived and felt.
\end{grammarlens}

\subsection*{Your turn}
\begin{itemize}
\raggedright
  \item \emph{Say it:} \emph{capṭī}
  \item \emph{Say it:} \emph{capṭī}, once more
  \item \emph{Say it:} \emph{śākmā̃ ek capṭī haḷdar nākho.} — Put a pinch of turmeric in the vegetables.
  \item \emph{Recall:} say the Gujarati for to temper, then the Gujarati for to sauté, then say \emph{capṭī} again
\end{itemize}

\begin{tcolorbox}[breakable,colback=teal!4,colframe=teal!35!black,title={Before you move on}]
What does \gu{ચપટી} mean? (\textquotedblleft{}A pinch (of something)\textquotedblright{}.) Say it once more.
\end{tcolorbox}
\section[\texorpdfstring{\gu{છાંટવું} (chā̃ṭvũ)}{છાંટવું (chā̃ṭvũ)}]{\gu{છાંટવું} (chā̃ṭvũ) — to sprinkle (a liquid)}

\label{lesson:GU-C520-chatvu}



\begin{quote}
Before the new one: say the Gujarati for to sauté, then the Gujarati for a pinch.
\end{quote}

\subsection*{You'll want to know: \gu{છાંટવું}}
\textbf{\gu{છાંટવું}} — \emph{chā̃ṭvũ} — \textquotedblleft{}to sprinkle (a liquid)\textquotedblright{}.

\begin{quote}
\gu{મીરાએ} \gu{ગરમ} \gu{તવા} \gu{પર} \gu{થોડું} \gu{પાણી} \gu{છાંટ્યું}.
\end{quote}

\begin{quote}
\emph{Mīrāe garam tavā par thoḍũ pāṇī chā̃ṭyũ.}
\end{quote}

\begin{quote}
Mira sprinkled a little water on the hot griddle.
\end{quote}

Past, with \gu{તેણે} (he or she …): \gu{તેણે} \gu{પાણી} \gu{છાંટ્યું}, \emph{teṇe pāṇī chā̃ṭyũ}.

\begin{cousinweb}
\textbf{\gu{છાંટવું}} (\emph{chā̃ṭvũ}) goes back to a reconstructed \emph{*chaṇṭ}, \textquotedblleft{}scatter\textquotedblright{}, and CDIAL lists the Gujarati verb, \textquotedblleft{}to sprinkle\textquotedblright{} (CDIAL 4970).
\end{cousinweb}

\begin{grammarlens}[title={What you've built}]
One more word for describing what you lived and felt.
\end{grammarlens}

\subsection*{Your turn}
\begin{itemize}
\raggedright
  \item \emph{Say it:} \emph{chā̃ṭvũ}
  \item \emph{Say it:} \emph{chā̃ṭvũ}, once more
  \item \emph{Say it:} \emph{Mīrāe garam tavā par thoḍũ pāṇī chā̃ṭyũ.} — Mira sprinkled a little water on the hot griddle.
  \item \emph{Recall:} say the Gujarati for to sauté, then the Gujarati for a pinch, then say \emph{chā̃ṭvũ} again
\end{itemize}

\begin{tcolorbox}[breakable,colback=teal!4,colframe=teal!35!black,title={Before you move on}]
What does \gu{છાંટવું} mean? (\textquotedblleft{}To sprinkle (a liquid)\textquotedblright{}.) Say it once more.
\end{tcolorbox}
\section[\texorpdfstring{\gu{ભૂંજવું} (bhū̃jvũ)}{ભૂંજવું (bhū̃jvũ)}]{\gu{ભૂંજવું} (bhū̃jvũ) — to parch (in a dry pan)}

\label{lesson:GU-C520-bhujvu}



\begin{quote}
Before the new one: say the Gujarati for a pinch, then the Gujarati for to sprinkle.
\end{quote}

\subsection*{You'll want to know: \gu{ભૂંજવું}}
\textbf{\gu{ભૂંજવું}} — \emph{bhū̃jvũ} — \textquotedblleft{}to parch (in a dry pan)\textquotedblright{}.

\begin{quote}
\gu{દાદીએ} \gu{તવા} \gu{પર} \gu{ચણા} \gu{ભૂંજ્યા}.
\end{quote}

\begin{quote}
\emph{dādīe tavā par caṇā bhū̃jyā.}
\end{quote}

\begin{quote}
Grandmother parched chickpeas on the griddle.
\end{quote}

Past, with \gu{તેણે} (he or she …): \gu{તેણે} \gu{ચણા} \gu{ભૂંજ્યા}, \emph{teṇe caṇā bhū̃jyā}.

\begin{cousinweb}
\textbf{\gu{ભૂંજવું}} (\emph{bhū̃jvũ}) comes from \emph{*bhṛñjati}, \textquotedblleft{}roasts\textquotedblright{}, a form CDIAL reconstructs, and CDIAL lists the Gujarati verb, \textquotedblleft{}to roast\textquotedblright{} (CDIAL 9586); it is said of grain roasted without oil.
\end{cousinweb}

\begin{grammarlens}[title={What you've built}]
One more word for describing what you lived and felt.
\end{grammarlens}

\subsection*{Your turn}
\begin{itemize}
\raggedright
  \item \emph{Say it:} \emph{bhū̃jvũ}
  \item \emph{Say it:} \emph{bhū̃jvũ}, once more
  \item \emph{Say it:} \emph{dādīe tavā par caṇā bhū̃jyā.} — Grandmother parched chickpeas on the griddle.
  \item \emph{Recall:} say the Gujarati for a pinch, then the Gujarati for to sprinkle, then say \emph{bhū̃jvũ} again
\end{itemize}

\begin{tcolorbox}[breakable,colback=teal!4,colframe=teal!35!black,title={Before you move on}]
What does \gu{ભૂંજવું} mean? (\textquotedblleft{}To parch (in a dry pan)\textquotedblright{}.) Say it once more.
\end{tcolorbox}
